\subsection{주요 결과}
\label{sec:results}


표~\ref{tab:main_results}에 제시된 결과는, 대부분의 파라미터를 과제 간 공유함에도 불구하고 PRAGMA가 평가된 거의 모든 도메인에서 기존 과제별 베이스라인을 일관되게 능가함을 보여준다.
가장 두드러진 향상은 영향이 큰 과제들의 정밀도-재현율 지표에서 관찰된다 — 신용 평가에서 PR-AUC가 130.2\,\% 증가했고, 커뮤니케이션 인게이지먼트에서 79.4\,\% 증가했다. 이는 PRAGMA가 전통 모델이 어려워하는 저빈도·고가치 신호를 식별하는 데 매우 효과적임을 시사한다.
ROC-AUC 이득은 더 점진적이지만, 같은 두 과제에서 각각 +12.4\,\%와 +20.4\,\%로 여전히 상당하다.
라이프타임 밸류와 정기 거래 같은 과제에서는 성능 차이가 더 좁혀지지만, 전체 추세는 PRAGMA가 고립된 과제별 모델의 성능과 맞먹거나 그 이상의 우수한 범용 표현을 제공함을 확인해 준다.

\begin{table}[t]
\centering
\begin{tabular}{llcr}
\textbf{과제}             & \textbf{지표}   & \textbf{베이스라인 (기준)} & \textbf{PRAGMA} \\
\midrule
신용 평가            & PR-AUC            & --                    & +130.2\,\%                \\
                          & ROC-AUC           & --                    &  +12.4\,\%                \\
\addlinespace
커뮤니케이션 인게이지먼트 & PR-AUC            & --                    &  +79.4\,\%                \\
                          & ROC-AUC           & --                    &  +20.4\,\%                \\
\addlinespace
외부 사기            & Precision         & --                    &  +16.7\,\%                \\
                          & Recall            & --                    &  +64.7\,\%                \\
\addlinespace
상품 추천    & mAP               & --                    &  +40.5\,\%                \\
\addlinespace
정기 거래    & $F_\text{1}$      & --                    &   +5.8\,\%                \\
\addlinespace
라이프타임 밸류            & PR-AUC            & --                    &   +1.8\,\%                \\
                          & ROC-AUC           & --                    &   +2.6\,\%                \\
\end{tabular}
\caption{
    \textbf{PRAGMA는 과제 간 대부분의 파라미터를 공유하면서도 내부 과제별 모델을 큰 폭으로 능가한다.}
    상대 성능은 ($\text{PRAGMA} / \text{baseline} - 1$)로 계산한다.
    PRAGMA는 LoRA 파인튜닝을 적용한 Large 변형을 사용한다.}
    \label{tab:main_results}
\end{table}


\subsubsection{모델 스케일의 효과}
\label{sec:scale_effect_result}
표~\ref{tab:scale_results}의 결과는, PRAGMA 아키텍처를 Small~(S, 10\,M) 변형에서 Medium~(M, 100\,M)과 Large~(L, 1\,B) 변형으로 확장했을 때의 성능 영향을 보여준다.
스케일링 이득은 과제별로 크게 다르며, 특히 신용 평가에 가장 크게 집중된다 — Large 모델은 Small 기준 대비 PR-AUC에서 +35.2\,\%, ROC-AUC에서 +5.8\,\% 이득을 낸다.

\begin{table}[t]
\centering
\begin{tabular}{llcrr}
& & \multicolumn{3}{c}{\textbf{PRAGMA}} \\
\cmidrule(lr){3-5}
\textbf{과제}             & \textbf{지표} & \textbf{S (기준)} & \textbf{M} & \textbf{L} \\
\midrule
외부 사기            & Precision       & --               &  +12.0\,\%             & +16.4\,\%             \\
                          & Recall          & --               &  +24.8\,\%             &  +23.5\,\%             \\
\addlinespace
상품 추천    & mAP             & --               &  +18.9\,\%             &  +27.0\,\%             \\
\addlinespace
신용 평가            & PR-AUC          & --               & +16.3\,\%             & +35.2\,\%             \\
                          & ROC-AUC         & --               &  +3.6\,\%             &  +5.8\,\%             \\
\addlinespace
라이프타임 밸류            & PR-AUC          & --               &  +1.5\,\%             &  +3.0\,\%             \\
                          & ROC-AUC         & --               &  +1.7\,\%             &  +3.4\,\%             \\
\addlinespace
커뮤니케이션 인게이지먼트  & PR-AUC          & --               & +0.1\,\%             & +1.6\,\%             \\
                          & ROC-AUC         & --               & $-$1.8\,\%             &  +0.7\,\%             \\
\addlinespace
정기 거래    & $F_\text{1}$    & --               &  +0.6\,\%             &  +0.4\,\%             \\
\end{tabular}
\caption{\textbf{모델 성능은 파라미터 수에 따라 스케일링된다}.
성능은 LoRA로 파인튜닝한 PRAGMA-S 대비이며, ($\text{model} / \text{PRAGMA-S} - 1$)로 계산한다.
}
\label{tab:scale_results}
\end{table}

특이하게도, 커뮤니케이션 인게이지먼트의 스케일링 거동은 단조롭지 않다 — Medium 변형은 ROC-AUC에서 약간의 회귀($-$1.8\,\%)를 보이고, Large 변형은 +0.7\,\%로 회복된다.
정기 거래나 LTV처럼 더 안정적인 지표에서는 성능 이득이 더 작아서, 보통 +3.5\,\% 이내에 머문다.
이 결과는 파라미터 수를 늘리면 일반적으로 예측력이 향상되지만, 거래성·라이프타임 밸류 예측에서는 Small 모델이 이미 매우 경쟁력 있는 표현을 제공한다는 것을 시사한다. 따라서 그러한 운영 사용 사례에 대해서는 Small이 잠재적인 효율성 스위트 스폿이 될 수 있다.


\subsubsection{사전학습의 효과}
\label{sec:pretraining_effect}

표~\ref{tab:scratch_vs_lora}의 결과는, 평가된 모든 과제에서 LoRA 파인튜닝이 처음부터의 전체 파라미터 학습 성능과 일관되게 맞먹거나 그 이상임을 보여 줌으로써 본 논문의 접근을 검증해 준다.
가장 큰 이득은 커뮤니케이션 인게이지먼트에서 관찰되는데, LoRA가 PR-AUC에서 +18.6\,\%, ROC-AUC에서 +5.0\,\%를 달성한다. 이는 사전학습된 PRAGMA 백본이, 단일 다운스트림 과제에서 처음부터 학습할 때는 배우기 힘든 풍부하고 다양한 이벤트 패턴을 포착하고 있음을 시사한다.
신용 평가도 비슷한 패턴을 따라, LoRA가 PR-AUC에서 +13.0\,\%, ROC-AUC에서 +1.6\,\%의 향상을 보인다.
상품 추천도 mAP에서 +10.3\,\%의 큰 이득을 본다.
정기 거래와 라이프타임 밸류에서는 향상폭이 더 작으며(각각 $F_1$ +0.6\,\%, PR-AUC / ROC-AUC +0.4\,\% / +0.3\,\%), 이는 처음부터 학습한 베이스라인이 이미 해당 목적에 관련된 구조를 대부분 포착했음을 의미한다. 이런 경우 LoRA 파인튜닝은 회귀 없이 동등한 성능을 유지한다.
이 결과들은 운영 환경에서 특히 의미가 크다. 즉, PRAGMA가 여러 독립적이고 유지보수 비용이 큰 모델들을 단일 공유 시스템으로 통합하면서도 예측 정확도를 희생하지 않고, 학습 가능한 파라미터 풋프린트는 훨씬 작게 유지할 수 있음을 확인해 주기 때문이다.

\begin{table}[t]
\centering
\begin{tabular}{llcr}
& & \multicolumn{2}{c}{\textbf{PRAGMA-M}} \\
\cmidrule(lr){3-4}
\textbf{과제} & \textbf{지표} & \textbf{Scratch (기준)} & \textbf{LoRA} \\
\midrule
커뮤니케이션 인게이지먼트 & PR-AUC          & --                     &  +18.6\,\%               \\
                         & ROC-AUC         & --                     &  +5.0\,\%               \\
\addlinespace
신용 평가           & PR-AUC          & --                     & +13.0\,\%               \\
                         & ROC-AUC         & --                     &  +1.6\,\%               \\
\addlinespace
상품 추천   & mAP             & --                     &  +10.3\,\%             \\
\addlinespace
정기 거래   & $F_\text{1}$    & --                     & +0.6\,\%              \\
\addlinespace
라이프타임 밸류           & PR-AUC          & --                     &  +0.4\,\%               \\
                         & ROC-AUC         & --                     &  +0.3\,\%               \\
\end{tabular}
\caption{\textbf{LoRA 파인튜닝과 처음부터 학습한 과제별 모델의 성능 비교.}
상대 성능은 ($\text{LoRA} / \text{Scratch} - 1$)로 계산한다.
LoRA는 처음부터의 전체 파라미터 학습 성능과 일관되게 맞먹거나 그 이상이다.}
\label{tab:scratch_vs_lora}
\end{table}
